\subsection{The cocharacter lattice and the actual quotient gauge}
\label{subsec:faithful-quotient-lattice-gauge}

Keep the field hypotheses and actual
\(\mathrm{AbQuot}(t,x)\) of
Proposition~\ref{prop:faithful-actual-abelian-quotient}.
Write \(G=G_x\), and write \(t_B\) for its proper polarized
quotient datum.  Choose actual germs \(g,g_B\) of dimensions
\(d,d'\).  We distinguish the actual quotient map \(q:G\to B\)
from its formal coordinate series \(Q_s\) and tangent matrix \(P\).

\begin{srproposition}[The actual \(O\)-lattice in the tangent sequence]
\label{prop:faithful-quotient-cocharacter-lattice}
There are a finite free abelian group \(\Lambda\), an actual
\(O\)-module structure on it, a surjective matrix
\(P:k^d\to k^{d'}\), and an injection
\[
                  j:k\otimes_{\mathbb Z}\Lambda\longrightarrow k^d
\]
with the following properties:
\[
 \begin{split}
 \ker P&=\operatorname{im}j,\\
 P R_a&=R_{B,a}P,\\
 j S_a&=R_a j,\qquad
 S_a=1\otimes(a|\Lambda).
 \end{split}
 \tag{QL1}\label{eq:faithful-quotient-lattice}
\]
The columns of \(P\) are precisely the actual first-order
coordinates of \(q\) in the given germs.
In particular no unrelated \(O\)-action is assigned to
\(\Lambda\) by a dimension comparison.
\end{srproposition}
\begin{proof}
Let \(C\) be the character group of the split torus in the
given extension, and put
\[
                         \Lambda=\operatorname{Hom}_{\mathbb Z}(C,\mathbb Z).
\]
It is finite free over \(\mathbb Z\).
The \(O\)-action on \(G\) preserves the torus kernel, since
\(q\) is equivariant for the actual action on \(B\).
Kernel factorization gives its induced group endomorphism
on \(D(C)\) for every \(a\in O\).
Every such endomorphism is a diagonal map defined by one
\(h(a):C\to C\): on the group algebra, the pullback of a
character is a group-like unit.  Comparing its Laurent
monomials, and then evaluating at the identity, makes it
a single character with coefficient \(1\).
Characters distinguish the torus maps, so uniqueness gives
the equations for the identity, multiplication, and addition.
This is the source's \srcid{exists\_diagMap} and
\srcid{exists\_charAct} construction.

Define the action on the cocharacter lattice by
\[
                       a\lambda=\lambda\circ h(a).
 \tag{QL2}\label{eq:faithful-cocharacter-action}
\]
Composition reverses the two character endomorphisms;
commutativity of \(O\) makes this a unital ring action, exactly
as in \srcid{lamEnd}.  It is not an action recovered merely
from its characteristic polynomial.

The tangent model of the split torus is
\(\operatorname{Hom}_{\mathbb Z}(C,k)\):
a tangent vector \(\delta\) sends a character \(\chi\) to
\(1+\varepsilon\delta(\chi)\) in the dual numbers.
The map
\[
 k\otimes_{\mathbb Z}\Lambda\longrightarrow
             \operatorname{Hom}_{\mathbb Z}(C,k),\qquad
 c\otimes\lambda\longmapsto
             \bigl(\chi\mapsto c\,\lambda(\chi)\bigr)
 \tag{QL3}\label{eq:faithful-cocharacter-tangent}
\]
is an isomorphism, by a basis of the finite free \(C\).
The torus action on this tangent model is precomposition
by \(h(a)\), which is exactly the base change of
\eqref{eq:faithful-cocharacter-action}.

The actual tangent sequence is
\[
       0\longrightarrow\operatorname{Lie}D(C)
        \stackrel{di}\longrightarrow\operatorname{Lie}G
        \stackrel{dq}\longrightarrow\operatorname{Lie}B
        \longrightarrow0 .
 \tag{QL4}\label{eq:faithful-actual-tangent-quotient}
\]
The source proves its first two assertions directly on
dual-number points.
The injection follows from the monomorphism \(i\).
If a tangent point of \(G\) maps to zero under \(q\), the
kernel pullback factors it uniquely through the torus.
Its reduction is the identity, again by injectivity of \(i\);
so this factor is a tangent point of the torus.
Conversely a torus tangent point maps to zero.
For surjectivity, \(q\) is smooth: after its flat surjective
quasi-compact base change by \(q\), it becomes the projection
\(D(C)\times G\to G\), which is smooth, and smoothness descends.
Given a dual-number point of \(B\) at the identity, smoothness
lifts it with the chosen identity reduction in \(G\).
Concretely the source factors the two points through their
local stalks, lifts the stalk algebra homomorphism over the
square-zero ideal of \(k[\varepsilon]/(\varepsilon^2)\), and
composes with the stalk-to-scheme map.
Thus it gives an actual tangent point and not only a
surjection after a field extension.

Choose the tangent-model linear isomorphisms
\(e_G:\operatorname{Lie}G\simeq k^d\) and
\(e_B:\operatorname{Lie}B\simeq k^{d'}\) determined by
the given germ bases.  Put
\[
             P=e_B\,dq\,e_G^{-1},\qquad
             j=e_G\,di\,(\text{map in \eqref{eq:faithful-cocharacter-tangent}}).
\]
Equation \eqref{eq:faithful-actual-tangent-quotient} gives
surjectivity, injectivity, and the kernel equality.
Differentiating the actual action equations for \(q\) and
\(i\), and using the character formula in
\eqref{eq:faithful-cocharacter-tangent}, gives the two
equivariance identities in \eqref{eq:faithful-quotient-lattice}.
Finally the germ basis condition identifies the image of
its \(l\)-th basis tangent point under \(q\) with the \(l\)-th
column of this same \(P\).  This is the actual
\(\mathrm{IsDiff}(g,g_B,P)\) field in \srcid{lie\_seq}.
\end{proof}

\begin{srlemma}[Axis-normalized units on the actual fibre square]
\label{lem:faithful-quotient-unit-axes}
If a global unit \(u\) on \(G\times G\) restricts to \(1\)
on both zero axes, then \(u=1\).
\end{srlemma}
\begin{proof}
The variety \(G\times G\) is integral: it is smooth and
geometrically connected over the algebraically closed field.
For the action of the torus on its first factor, compare
the translated \(u\) with \(u\).
The same Laurent-unit coefficient argument as in
Lemma~\ref{lem:faithful-torus-quotient-seesaw} says that
the ratio is a fixed character of the torus.
Evaluate at a point of the second zero axis:
both the original and translated values are \(1\).
This makes that character trivial.
Do the same for the torus action in the second factor,
using the first zero axis.
Thus \(u\) is invariant under the two torus actions.

The map \(q\times q:G\times G\to B\times B\) is flat,
surjective, and quasi-compact.  Two points in a fibre
differ by these two kernel-torus translations.
The two pullbacks of \(u\) to its fpqc relation agree,
so \(u\) descends to a global function on \(B\times B\).
That proper geometrically connected smooth variety has
global functions \(k\).  The descended function is constant,
and a zero-axis value makes its constant \(1\).
This is the actual \srcid{UnitTrivialOnBothAxes} proof.
\end{proof}

\begin{srproposition}[The actual mixed-coefficient gauge]
\label{prop:faithful-quotient-mixed-gauge}
There is one matrix \(Z\in\operatorname{Mat}_d(k)\) such that,
for every \(a\in O\),
\[
          C_a+R_a^t Z-ZR_{\bar a}=P^t C_{B,a}P .
 \tag{QL5}\label{eq:faithful-quotient-mixed-gauge}
\]
It is obtained from the actual line isomorphism
\(L_x\simeq q^*P_B\) and the two chosen actual formal frames.
\end{srproposition}
\begin{proof}
We follow the finite-jet construction and comparison in
\srcid{AbelianQuotientBilinear}, retaining its unit normalization.

First the universal coordinates of \(g_B\) express
\(q\circ g\) at every finite jet order by a unique truncated
coordinate tuple.  Truncating one such expression gives the
previous one by uniqueness.  Its constant term is zero,
since \(q(0)=0\).
The inverse limit of the compatible coefficients gives
a tuple of formal series
\[
                     Q_s\in k[[X_1,\ldots,X_d]]^{d'},
                      \qquad Q_s(0)=0 .
\]
At every order it represents the actual point \(q\circ g\).
The first-order uniqueness for the actual germ bases, and
the last assertion of
Proposition~\ref{prop:faithful-quotient-cocharacter-lattice},
give \(D Q_s(0)=P\).
Applying the actual equivariance of \(q\) to the universal
jet points and using coordinate uniqueness also gives the
formal action compatibility
\[
             Q_s([a]_gX)=[a]_{g_B}(Q_s(X)).
 \tag{QL6}\label{eq:faithful-quotient-formal-action}
\]
These are the source's
\srcid{AbelianQuotientGermPowerSeries},
\srcid{AbelianQuotientDifferentialMatrix}, and
\srcid{AbelianQuotientGermSubstitution} steps.

Let \(w:L_x\simeq q^*P_B\) be the actual line isomorphism.
On the two universal jet points, it induces an isomorphism
of their normalized second-difference lines.
The chosen frame \(e_G\) and the pullback of \(e_B\) are
generators of these identified lines, so their ratio is a
unique unit \(U_m\) in the finite two-point jet algebra.
Use the source's orientation:
\[
               e_G=U_m^{-1}\,\gamma_w\,Q_s^*e_B ,
 \tag{QL7}\label{eq:faithful-quotient-frame-ratio}
\]
where \(\gamma_w\) is the alternating product induced by
\(w\) on the second difference.
The source records this as
\(\mathrm{Tr}(w,e_G,Q_s^*e_B,U_m^{-1})\).
Restriction of the two actual frames and of \(w\) shows
that \(U_{m+1}\) truncates to \(U_m\); equality follows
by the uniqueness of the ratio of generators.
Thus these units give a formal series
\[
                   \lambda(X,Y),\qquad
                   \lambda(X,0)=\lambda(0,Y)=1 .
\]
For the two axis identities, restrict
\eqref{eq:faithful-quotient-frame-ratio} to that axis.
Both normalized second-difference generators become the
canonical axis generator, and the alternating product of
\(w\) cancels in pairs.  The resulting ratio is \(1\)
at every jet order.  This is the precise
\srcid{AbelianQuotientGaugeTransform.exists\_lam}
construction, including its inverse orientation.

For each \(a\), compare the two actual global adjoint
generators \(w_{G,a}\) and \(w_{B,a}\) using \(w\).
Their alternating line expressions coincide because \(q\)
is a homomorphism and is equivariant for both \(a\) and
\(\bar a\).  Their ratio is therefore a global unit
\(u_a\) on \(G\times G\);
this is \srcid{AbelianQuotientLineBundleGauge}.
The source checks that it restricts to \(1\) on both axes.
It first restricts to the actual universal one-point jets.
The adjoint relation makes both generator values \(1\);
the alternating product of \(w\) along an axis again
cancels in pairs, so the pullback of \(u_a\) is \(1\)
at every order.
These jet equalities imply the actual axis equality:
the smooth germ coordinates identify all truncated local
rings, their intersection is zero by Krull intersection,
and the restriction of a global function on the integral
\(G\) to its local ring at the identity is injective.
No surjectivity of the formal Taylor map is used.
This is \srcid{AbelianQuotientGaugeUnit}.
Lemma~\ref{lem:faithful-quotient-unit-axes} now gives
\(u_a=1\) on the whole fibre square.

Put \(A_a(X,Y)=g.B.\mathrm{ad}(a)\) and
\(A_{B,a}=g_B.B.\mathrm{ad}(a)\).
The two adjoint-frame equations and
\eqref{eq:faithful-quotient-frame-ratio}, with this global
ratio equal to \(1\), give at every finite jet order
\[
 \begin{split}
 A_{B,a}(Q_s(X),Q_s(Y))\,
                    \lambda(X,[\bar a]_gY)
     =\lambda([a]_gX,Y)\,A_a(X,Y).
 \end{split}
 \tag{QL8}\label{eq:faithful-quotient-series-gauge}
\]
For clarity, the cancellation uses the same four factors
at each order: the two adjoint-frame equations
\(w_{G,a}e_{G,\mathrm{left}}=A_a e_{G,\mathrm{right}}\)
and
\(w_{B,a}e_{B,\mathrm{left}}=A_{B,a}e_{B,\mathrm{right}}\),
and \eqref{eq:faithful-quotient-frame-ratio} at the left
and right substituted point pairs.
The product induced by \(w\) on the adjoint line is the
right second-difference product divided by the left one;
its factors cancel because \(u_a=1\).
Equation \eqref{eq:faithful-quotient-formal-action} ensures
that substituting \(Q_s\) before or after either action
gives the same truncated point.
The remaining unit identity is exactly
\eqref{eq:faithful-quotient-series-gauge}.
Equality at every finite order gives equality of formal
series, as in \srcid{AbelianQuotientGaugeRelation}.

Set \(Z=\operatorname{bil}_{1,1}\lambda\).
All four unit series in
\eqref{eq:faithful-quotient-series-gauge} are \(1\)
on both axes.  Therefore the mixed coefficient of a product
is the sum of the mixed coefficients.
A substitution contributes only its first differential:
\[
 \begin{split}
 \operatorname{bil}\bigl(A_{B,a}(Q_s(X),Q_s(Y))\bigr)
                                  &=P^t C_{B,a}P,\\
 \operatorname{bil}\bigl(\lambda(X,[\bar a]_gY)\bigr)
                                  &=Z R_{\bar a},\\
 \operatorname{bil}\bigl(\lambda([a]_gX,Y)\bigr)
                                  &=R_a^t Z.
 \end{split}
\]
Taking mixed coefficients in
\eqref{eq:faithful-quotient-series-gauge} gives
\[
                    P^t C_{B,a}P+Z R_{\bar a}=R_a^t Z+C_a,
\]
which is \eqref{eq:faithful-quotient-mixed-gauge}.
This is one \(Z\), obtained from the single ratio series
\(\lambda\), for all \(a\).
\end{proof}

\paragraph{Frozen-source ledger and scope.}
The actual lattice construction is
\href{https://github.com/ChristianKudeba/alpoge-rh-verification/blob/alpoge-source-7cbf7a0c/src/Development/AutomorphicSpectrum/Proofs/GeometryProofs/AbelianQuotientLattice.lean}
{\srcid{AbelianQuotientLattice}}
(blob \srcid{187d2323635424a4d4fc7f2e08fd28c423478a39}):
\srcid{exists\_dual\_lift}, \srcid{lie\_kernel},
\srcid{lie\_surjective}, \srcid{exists\_charAct},
\srcid{tenMap\_bijective}, \srcid{lamEnd}, and
\srcid{lie\_seq}.
The mixed-gauge assembly is
\href{https://github.com/ChristianKudeba/alpoge-rh-verification/blob/alpoge-source-7cbf7a0c/src/Development/AutomorphicSpectrum/Proofs/GeometryProofs/AbelianQuotientBilinear.lean}
{\srcid{AbelianQuotientBilinear.frame\_gauge}}
(blob \srcid{c5ea956019648f4e8862164cedc13627b0d6246e}).
Its actual finite-jet producers, in that same source directory, are
\srcid{AbelianQuotientGermPowerSeries}
(\srcid{b8393721074d0e17b4eb268b369cdcd4f32c8f7e}),
\srcid{AbelianQuotientDifferentialMatrix}
(\srcid{8bac6126d6aaf9d18aa282463b4bc805bc315bc5}),
\srcid{AbelianQuotientGermSubstitution}
(\srcid{7781ab82d410a2d0061107d0a64f4823a7241c88}),
\srcid{AbelianQuotientGaugeTransform}
(\srcid{c997b24547de07406ad472c9b9e9846cb88b36ee}),
\srcid{AbelianQuotientLineBundleGauge}
(\srcid{c27e04cd714a1cff5fe6806b1347cf0b94dc1f3b}),
\srcid{AbelianQuotientGaugeUnit}
(\srcid{3e59e5015dc988a8a5a481a93a11364145b4f5b4}),
\srcid{UnitTrivialOnBothAxes}
(\srcid{196eb507e5f9a50a5d4a8ea317ca2cb707deab37}),
and \srcid{AbelianQuotientGaugeRelation}
(\srcid{5352d8f5464c4871367b864a9724367ec83dbb75}).
The assembly, lattice, and eight displayed helper files were checked
against the frozen baseline.
The bilinear coefficient step is the actual
\srcid{FormalBilinearTermProduct} and
\srcid{SubstitutionBilinearFormTransform} calculation.

The propositions construct the actual lattice, action,
differential, and gauge on the supplied geometric quotient.
They do not infer an integral formal-Hodge identification
from its dimension, nor do they assert that a nonflat
base change preserves an eigensaturated lattice.
The remaining geometric input to
\srcid{HeightCharFreeGeometric} is the finite-kernel
norm-degree producer for the proper quotient.
